%==============================================================================
%  08-case-studies.tex — Section 7: Benchmark Case Studies
%==============================================================================
\strategyPart{Section 07 \textbar\ Case Studies}{Benchmark Case Studies and Research Insights}

Two organizations — one in financial services, one in consumer goods — show
how the problem-first philosophy produces reinvention at scale. In both cases,
the AI was not the strategy; the organizational problem was.

\section{Morgan Stanley: reinventing the advisor--client interaction}

Morgan Stanley identified an expensive organizational problem: its
\~{}50,000 financial advisors could not search the firm's vast research
library quickly enough to answer client questions in real time. In 2023 the
firm launched \textit{AI @ Morgan Stanley Assistant}, built on OpenAI's
GPT-4, giving advisors a conversational interface to internal research and
content [Business Insider 2023; Morgan Stanley 2023]. This is a Level 1
(Assist) deployment with a precise problem definition and a measurable
outcome: advisor productivity and client responsiveness. The lesson is not
``Morgan Stanley adopted GPT-4''; it is that a knowledge-retrieval bottleneck
in a high-value workflow was solved with the smallest sufficient capability.

\section{Coca-Cola: creative reinvention and optimization at scale}

Coca-Cola attacks two distinct problems. On the creative side, its
\textit{Create Real Magic} generative-AI platform — built with OpenAI and
Bain — lets marketers and creators produce on-brand content at scale,
attacking a content-production bottleneck across a global brand portfolio
[Forbes 2023]. On the operational side, the company has long used predictive
and optimization AI for demand forecasting and supply-chain planning, and in
2024 committed over \$1~billion to a five-year cloud and generative-AI
partnership with Microsoft [Reuters 2024]. The pattern is textbook
problem-first portfolio construction: generative AI for the knowledge and
creativity problem, optimization AI for the resource-allocation problem.

\section{Synthesis: problems beat tools}

The contrast across the research base is consistent. Organizations that frame
AI around business problems and KPIs reach production and scale; organizations
that frame AI around technology adoption stall at proof of concept — the exact
population captured by Gartner's abandonment prediction [Gartner 2023].
McKinsey's economic analysis reinforces the point: roughly three-quarters of
generative-AI value sits in a handful of business functions where the problems
are biggest [McKinsey 2023]. The case-study evidence and the macro evidence
point to the same conclusion: \textbf{the organization that starts with the
problem wins; the organization that starts with the tool waits.}
